\Qhead{1.3}{Воздействие гармонического сигнала на автоколебательные системы}{Воздействие гармонического сигнала на автоколебательные системы. Синхронизация. Явления захватывания и гашения колебаний. Применение захватывания для стабилизации частоты.}
\Subfile{1.3.1}{sections/01_oscillations/q03/s01}
\Subfile{1.3.2}{sections/01_oscillations/q03/s02}
\Subfile{1.3.3}{sections/01_oscillations/q03/s03}

\recap{Свободная фаза позволяет слабому сигналу захватить автогенератор. Уравнение Адлера даёт полосу и биения. Сильный сигнал может погасить самостоятельную составляющую через насыщение; захват стабильным источником используется для стабилизации частоты.}
\board{Уравнение Адлера; «галочка» частоты биений с нулевым отрезком; механизм насыщения активного элемента.}
\oral{Почему одна линия на выходе ещё не позволяет различить захват и гашение?}
\bait{Можно оставить замедление биений у края захвата: фаза долго проходит участок почти компенсированной расстройки.}
\sources{Пятаков~\cite{pyatakov}, PDF-с.~63--65: внешнее воздействие, захват и тушение; вопрос~3 коллеги --- основа объяснения. Уравнения Адлера, размерная нормировка силы и двухчастотное усреднение воспроизведены здесь; непроверенные коэффициенты конкретного прибора не заимствованы.}
